\chapter{Einführung}
\index{Agentic AI}%
\index{Software Engineering}%

% Längeres abgesetztes Zitat
\begin{quote}
\textit{\textquote[\cite{karpathy2023intro}]{Wir erleben die Entstehung eines neuen Betriebssystems: Das LLM fungiert als Kernel-Prozess, das Kontextfenster als Arbeitsspeicher und externe Werkzeuge als E/A-Schnittstellen. Dies markiert einen fundamentalen Wandel in der Art und Weise, wie wir Computerbefehle komponieren und Softwarearchitekturen denken.}}
\smallskip
\begin{flushright}
— Andrej Karpathy, \textit{Intro to Large Language Models}, 2023
\end{flushright}
\end{quote}

\section{Motivation und Relevanz}

Im Software Engineering zeichnet sich derzeit eine Verschiebung der Automatisierungslogik ab: \emph{agentische KI} \textendash{} also KI-Systeme, die Ziele interpretieren, Pläne erstellen, Werkzeuge verwenden und Ergebnisse eigenständig prüfen \textendash{} erweitert klassische Automatisierung um adaptive, mehrschrittige Problemlösung.\footnote{Der Begriff \emph{agentische KI} bezeichnet Systeme, die über mehrere Schritte hinweg komplexe Aufgaben bearbeiten und dabei Planung, Werkzeugnutzung und Rückkopplung kombinieren.}
 
Darüber hinaus verlangt die praktische Anwendung agentischer Systeme einen belastbaren Ausgleich zwischen Automatisierung und menschlicher Kontrolle. In vielen Entwicklungsprozessen sind Schritte, die bislang überwiegend manuell erfolgen \textendash{} etwa Code-Reviews, Testgenerierung oder Release-Checks \textendash{} geeignete Kandidaten für assistierte Abläufe. Agenten können Routinearbeit übernehmen, erzeugen damit jedoch zugleich neue Anforderungen an Validierung, Protokollierung und Eingriffsmöglichkeiten. Die vorliegende Arbeit untersucht diesen Zielkonflikt und entwickelt Gestaltungsprinzipien für einen kontrollierten produktiven Einsatz.
Wie Richards und Ford betonen, \textquote[\cite{richards2020fundamentals}]{müssen Architekturprinzipien auf Skalierbarkeit und Wartbarkeit ausgelegt sein}, um praktischen Anforderungen gerecht zu werden.
Für Forschung und Praxis ergeben sich daraus neue Architektur- und Methodikfragen: Wie lassen sich robuste agentische Workflows entwerfen? Wie werden Werkzeugnutzung, Gedächtnis und Langkontext so orchestriert, dass nachvollziehbare Entscheidungen entstehen? Und wie lassen sich Sicherheit, Prüfbarkeit und Testbarkeit systematisch in agentische Systeme integrieren?\footnote{Praktische Orchestrierung bezeichnet hier die Koordination von Planungsschritten, Werkzeugaufrufen und Gedächtniszugriffen in einer strukturierten Abfolge.}

\section{Problemstellung}

Vor dem Hintergrund adressiert die Arbeit exemplarisch die Entwicklung und Evaluation eines agentischen Systems für Softwareentwicklungsaufgaben (z.\,B. Refactoring, Code-Review, Generierung von Tests). Zentrale Fragen sind:

\begin{itemize}
  \item Wie lassen sich Agentenstrategien (Planen, Werkzeugaufrufe, Selbstkritik) systematisch modellieren?
  \item Wie werden externe Werkzeuge (VCS, CI, Linter, Issue-Tracker) sicher und nachvollziehbar eingebunden?
  \item Welche Metriken messen Fortschritt, Qualität und Sicherheit realistisch?
\end{itemize}
 
Die hier formulierten Probleme sind sowohl konzeptioneller als auch praktischer Natur: Im Mittelpunkt stehen nicht nur die abstrakte Modellierung von Strategien, sondern ebenso konkrete Implementierungsfragen (Schnittstellen, Serialisierung, Fehlerbehandlung) und Fragen des Evaluationsdesigns (Benchmarks, Messgrößen, Reproduzierbarkeit). Die Ergebnisse sollen Entwicklungsteams konkrete Handlungsempfehlungen geben, wie agentische Automatisierung schrittweise, kontrolliert und messtechnisch abgesichert eingeführt werden kann.\index{Evaluation}
\section{Zielsetzung}

Die Ziele der Arbeit sind auf die Spezialisierung \enquote{Software Engineering mit agentischer KI} zugeschnitten:\footnote{Siehe auch~\cite{wang2023survey} für verwandte Arbeiten zu Agentenarchitekturen.}

\begin{enumerate}
  \item Analyse des Forschungsstands zu agentischen Architekturen und Orchestrierungsframeworks
  \item Entwurf einer referenzierbaren Agentenarchitektur für Software-Engineering-Aufgaben
  \item Implementierung eines prototypischen Agenten mit Werkzeuganbindung und Gedächtnis
  \item Evaluation anhand reproduzierbarer Benchmarks (Qualität, Kosten, Laufzeit, Sicherheit)
\end{enumerate}
 
Kurz gefasst zielt die Arbeit darauf ab, aus Praxisproblemen generalisierbare Lösungen abzuleiten: Neben einem lauffähigen Prototyp steht die Frage im Vordergrund, welche Architekturmuster sich zuverlässig übertragen lassen und welche Messgrößen den Mehrwert gegenüber manuellen Prozessen belastbar quantifizieren. Damit richtet sich die Arbeit gleichermaßen an Forschende und an Verantwortliche in der Softwareentwicklung.
\section{Abgrenzung des Themas}

Die Arbeit fokussiert Agenten für Softwareentwicklungsaufgaben. Nicht im Fokus sind z.\,B. Reinforcement Learning from Human Feedback (RLHF) im Detail, Trainingsmethoden auf Rohdaten oder proprietäre Interna von Foundation Models. Ebenso werden Domänen außerhalb der Softwareentwicklung (z.\,B. Robotik) nicht betrachtet.

\begin{itemize}
  \item Zu komplexe Spezialfälle, die für die Arbeit nicht relevant sind
  \item Historische Entwicklungen vor einem bestimmten Zeitpunkt
  \item Randbereiche, die außerhalb des Fokus liegen
\end{itemize}
 
Die Konzentration auf reproduzierbare Ergebnisse erlaubt es, bewusst auf tiefe Trainingsfragen zu verzichten; stattdessen wird die Interaktion mit bestehenden Foundation Models über APIs und Adapterlösungen untersucht. Diese Fokussierung erhöht die Nachvollziehbarkeit der Ergebnisse und verbessert die Anschlussfähigkeit an typische Software-Engineering-Umgebungen.
\section{Aufbau der Arbeit}

Der Aufbau der Arbeit folgt einem klaren Argumentationsgang: Zunächst werden in Kapitel~\ref{ch:hintergrund} die relevanten Grundlagen und verwandten Arbeiten zusammengetragen. Darauf aufbauend beschreibt Kapitel~\ref{ch:konzept} die entworfene Referenzarchitektur mit Fokus auf Zustandsmodellierung, Policy-Design und Schnittstellen. Kapitel~\ref{ch:realisierung} dokumentiert die Implementierung des Prototyps, zeigt zentrale Code-Beispiele und erläutert Integrationsdetails. Abschließend fasst Kapitel~\ref{ch:abschluss} die Ergebnisse zusammen, diskutiert Limitationen und gibt einen Ausblick auf weiterführende Forschungs- und Entwicklungsfragen.

Insgesamt soll die Arbeit nicht nur eine theoretische Diskussion bieten, sondern konkrete, übertragbare Architekturentscheidungen und Metriken bereitstellen, die in produktiven Entwicklungsumgebungen einsetzbar sind.\index{Agentenarchitektur}
